% LLNCS/Springer draft for SOICT-style proceedings.
\documentclass[runningheads]{llncs}

\usepackage[T1]{fontenc}
\usepackage{lmodern} % Vector CM-compatible fonts; this host lacks CM-Super Type 1.
\usepackage{graphicx}
\usepackage{amsmath}
\usepackage{amssymb}
\usepackage{float}

\begin{document}

\title{MixiMotion: A One-Step HY-Motion-1.0 Model for Text-to-Motion Generation}
\titlerunning{MixiMotion}

\author{Anonymous Authors}
\authorrunning{Anonymous Authors}
\institute{Anonymous Institution\\
\email{anonymous@example.com}}

\maketitle

\begin{abstract}
Diffusion and flow-matching text-to-motion models can generate high quality
3D human motion, but their multi-step sampling makes interactive deployment
expensive. We present MixiMotion, our compact one-step HY-Motion-1.0 model
designed to preserve semantic alignment while reducing inference to a single
network evaluation. MixiMotion is trained with a
progressive curriculum that first learns a one-step motion prior, then adds
motion-quality curricula for contact and locomotion, followed by semantic
balancing and preference optimization. Training uses cached motion samples
produced by HY-Motion-1.0-Lite. Before retraining, we remove every cached row
whose normalized prompt exactly matches an official Structured Semantic
Alignment Evaluation (SSAE) prompt. On 1,969 videos per system, an all-frame
native-video evaluation with Qwen/Qwen3-VL-30B-A3B-Instruct-FP8 gives 0.8802
for MixiMotion, compared with 0.8856 for HY-Motion-1.0 full, 0.8764 for
MotionHiFlow, and 0.8415 for MotionLCM. Under this judge, MixiMotion retains
99.4\% of the full model's SSAE score while using 460.0M rather than 1.04B
motion-generator parameters and requiring one network evaluation per motion
candidate. A blinded human study with 25 participants further rates
MixiMotion at 4.33/5 for overall motion quality, compared with 4.50 for
HY-Motion-1.0 full, 4.20 for MotionHiFlow, and 3.98 for MotionLCM.

\keywords{Text-to-motion \and Knowledge distillation \and One-step generation \and Human motion synthesis \and Semantic alignment}
\end{abstract}

\section{Introduction}

Text-to-motion generation converts natural-language descriptions into 3D human
motion sequences. Recent flow-matching and diffusion-transformer systems such
as HY-Motion-1.0~\cite{hymotion2025} provide strong semantic alignment and natural
motion, but their iterative solvers are costly for interactive authoring,
game engines, and web demos. A natural alternative is to distill the teacher
into a one-step generator. However, naively matching teacher latents is not
sufficient: motion quality constraints such as foot contact, root drift, and
locomotion stability interact with semantic behaviors involving hands, torso,
and object-like actions.

This paper summarizes MixiMotion, our current one-step HY-Motion-1.0
text-to-motion system. Rather than presenting a sequence of implementation
versions, we frame the training process as a
progressive curriculum. The curriculum starts from teacher imitation, adds
motion-quality objectives, introduces balanced semantic groups, and finishes
with preference optimization over teacher-generated candidates. This framing is
important for reproducibility: the individual checkpoints are implementation
milestones, whereas the method is the loss and data schedule.

Our current experimental focus is not to claim a new state of the art over
HY-Motion-1.0 full, but to quantify how close the one-step MixiMotion model
can get to the full multi-step baseline under the same VLM-based semantic
evaluation. The primary HY-Motion-family comparison shares the renderer and
video settings; public baselines use the same native-video judge and questions
but retain their native renderers.

Our contributions are threefold:
\begin{itemize}
\item We formulate one-step HY-Motion distillation as bidirectional set
matching between multiple student noise samples and multiple cached teacher
motions, with an objective designed to preserve multi-sample coverage in a
single network evaluation.
\item We study a progressive training recipe that augments latent imitation
with decoded motion constraints, semantic-group balancing, and
distance-based preference fine-tuning. Our ablation reports non-monotonic
outcomes rather than assuming that every stage is beneficial.
\item We evaluate with native-video SSAE after removing exact text matches
from the training bank. We report paired bootstrap intervals and a
common-hardware efficiency comparison against HY-Motion-1.0 full,
MotionHiFlow, and MotionLCM, together with a blinded human evaluation of
semantic alignment, naturalness, and overall motion quality.
\end{itemize}

Figure~\ref{fig:overview} summarizes the offline teacher-cache construction and
the one-step student supervision used to train MixiMotion.

\begin{figure}[t]
\centering
\includegraphics[width=\textwidth]{figures/miximotion_overview.png}
\caption{Overview of MixiMotion distillation. A frozen HY-Motion-1.0-Lite
teacher is run once to cache motion targets and text features. MixiMotion maps
Gaussian noise to one-step motion candidates, supervised by latent set matching
and decoded-motion matching; gradients update only the student.}
\label{fig:overview}
\end{figure}

\section{Related Work}

\paragraph{Text-to-motion generation.}
Diffusion and flow-based models have substantially improved the diversity and
semantic fidelity of text-conditioned human motion. HumanML3D established a
large text-motion benchmark~\cite{humanml3d2022}; T2M-GPT models discrete
motion tokens autoregressively~\cite{t2mgpt2023}; and MLD performs diffusion in
a compact motion latent space~\cite{mld2023}. MDM~\cite{mdm2022} introduced a
strong motion-space diffusion formulation, while HY-Motion-1.0 scales
flow-matching transformers and semantic evaluation~\cite{hymotion2025}.
Efficiency-oriented work includes latent consistency training in
MotionLCM~\cite{motionlcm2024}, one-step flow matching in
MoFlow~\cite{moflow2024}, and hierarchical flow matching in
MotionHiFlow~\cite{motionhiflow2026}. MixiMotion differs in targeting the
HY-Motion latent and conditioning interface directly, allowing a controlled
comparison with the full HY-Motion-1.0 generator.

\paragraph{Fast generative-model distillation.}
Progressive distillation reduces iterative samplers by repeatedly teaching a
student to replace multiple teacher steps~\cite{progressivedistill2022}, while
consistency models learn outputs that agree across points of a diffusion
trajectory~\cite{consistencymodels2023}. Our setting instead caches several
complete teacher motions per prompt and learns a direct one-pass mapping. This
avoids online teacher evaluation but requires the student to cover a
multi-modal target set.

\paragraph{Set matching and preference learning.}
Implicit maximum likelihood estimation matches generated samples to nearby
observations and offers an alternative to pointwise regression for
multi-modal generation~\cite{imle2018}. We use bidirectional nearest-neighbor
coverage over student and teacher sets. In the final stage, we adapt the
relative-preference principle of DPO~\cite{dpo2023} to distance-based motion
scores with a frozen pre-preference student as reference. Unlike reward-model
training, preference pairs are derived from deterministic physical statistics
of cached teacher candidates.

\section{Method}

\subsection{Teacher and MixiMotion}

Let $G_T$ denote a multi-step HY-Motion-1.0 teacher and $G_\theta$ denote
MixiMotion. Given a text prompt $p$, the teacher produces a set of candidate
normalized motion latents
\begin{equation}
  \mathcal{Y}_T(p) = \{y_1, \ldots, y_K\},
\end{equation}
where each candidate is generated with a deterministic seed. MixiMotion
receives the same text representation and predicts a motion latent in one
network evaluation. Concretely, the student input consists of Gaussian noise
tokens $z \in \mathbb{R}^{T \times D}$, token-level text context $C(p)$, a
global text vector $v(p)$, and a temporal validity mask $m$. The student maps
\begin{equation}
  \hat{y} = G_\theta(z, C(p), v(p), m), \qquad z \sim \mathcal{N}(0, I),
\end{equation}
where $T$ is the motion length and $D$ is the latent dimension. The implemented
network encodes the noise tokens, text context, and text vector separately,
forms a conditioning adapter from a masked noise summary and the text vector,
passes motion and text features through double-stream and single-stream
multimodal DiT blocks, and predicts the final normalized motion latent. A
frozen teacher decoder is used only to recover joint trajectories for
motion-quality analysis and video rendering.

\subsection{Progressive Curriculum Distillation}

The training recipe has four stages. Let
$\mathcal{X}_\theta(p)=\{\hat{y}_1,\ldots,\hat{y}_M\}$ be $M$ student
candidates generated from independent Gaussian noise tokens for the same
prompt. For two motion latents $a,b$ with valid length $\ell$, we use a masked
Smooth-L1 distance
\begin{equation}
  d_\delta(a,b;\ell)
  = \frac{1}{\ell D} \sum_{t=1}^{\ell} \sum_{d=1}^{D}
  \rho_\delta(a_{t,d}-b_{t,d}),
\end{equation}
where $\rho_\delta$ is the Smooth-L1 penalty with transition point $\delta$.

\subsubsection{Base one-step distillation.}
The first stage learns the teacher distribution using a set-to-set IMLE-style
objective. For each prompt, MixiMotion draws several candidates and is matched
to the closest teacher candidate. Unlike a single nearest-neighbor loss, the
implemented objective uses both teacher-to-student and student-to-teacher
coverage terms:
\begin{equation}
  \mathcal{L}_{\mathrm{set}}
  = \alpha_{\mathrm{t2s}}\frac{1}{K}\sum_{j=1}^{K}
    \min_{i \le M} d_\delta(\hat{y}_i,y_j;\ell)
  + \alpha_{\mathrm{s2t}}\frac{1}{M}\sum_{i=1}^{M}
    \min_{j \le K} d_\delta(\hat{y}_i,y_j;\ell).
\end{equation}
The implemented base loss also includes endpoint and trajectory terms in the
normalized representation and decoded body-space supervision:
\begin{equation}
  \mathcal{L}_{\mathrm{base}}
  = \mathcal{L}_{\mathrm{set}}
  + \lambda_{\mathrm{end}}\mathcal{L}_{\mathrm{end}}
  + \lambda_{\mathrm{traj}}\mathcal{L}_{\mathrm{traj}}
  + \sum_{r\in\mathcal{R}_{\mathrm{base}}}\lambda_r\mathcal{L}_r .
\end{equation}
Here, $\mathcal{R}_{\mathrm{base}}$ contains decoded body, foot, key-joint, and
contact terms. These terms use the same temporal mask and bidirectional
nearest-match reduction where a candidate set is involved. This stage gives
MixiMotion a broad one-step motion prior while exposing the latent predictor
to trajectory and body-space errors that latent distance alone may hide.

\subsubsection{Motion-quality curriculum.}
The second stage decodes both teacher and student candidates and applies
bidirectional nearest-neighbor matching in body space. The supervision covers
body position, velocity, and acceleration; foot position, velocity, and
contact velocity; stride, foot lift, and alternating contact; and root, yaw,
and arm-motion statistics. Training first performs gait/body matching, then a
hard-locomotion subphase oversamples locomotion prompts by a factor of six and
increases the emphasis on stride, lift, and contact. These terms target
sliding, weak steps, and unstable lower-body motion without introducing an
online teacher.

\subsubsection{Semantic-balanced curriculum.}
The third stage balances two lexical training lanes: locomotion and
stationary/localized actions. Prompts in both lanes are oversampled by a
factor of four. Locomotion examples retain stride, foot-lift, and alternating
contact supervision, whereas localized actions such as reaching, waving,
sitting, kneeling, or squatting receive root-displacement and root-velocity
penalties to discourage unintended translation. The anchored and other groups
are not separate lanes in this stage; they are introduced only by the
group-aware preference scoring described next.

\subsubsection{Preference optimization.}
Finally, we build preference pairs from teacher candidates. For each prompt,
the best and worst teacher candidates are selected according to the
group-aware motion-quality score $q(y,p)$. For locomotion prompts, $q$ rewards
foot lift, alternating contact, and stride while penalizing skating and contact
velocity. For stationary and anchored prompts, $q$ penalizes root displacement,
root velocity, skating, and unnecessary foot lift. Other prompts use a mixed
score that mildly rewards visible motion while penalizing instability. The
winner and loser are
\begin{equation}
  y^+ = \arg\max_{y_j \in \mathcal{Y}_T(p)} q(y_j,p),
  \qquad
  y^- = \arg\min_{y_j \in \mathcal{Y}_T(p)} q(y_j,p).
\end{equation}
MixiMotion is then optimized with a DPO-style preference objective plus
supervised fine-tuning toward the preferred candidate. A frozen reference
student $G_{\mathrm{ref}}$ is initialized from the model before preference
training and evaluated on the same sampled noise tokens. Define the best-of-$M$
score
\begin{equation}
  s_\theta(y) = -\min_{i \le M} d_\delta(\hat{y}_i,y;\ell),
  \qquad
  s_{\mathrm{ref}}(y) = -\min_{i \le M}
  d_\delta(\hat{y}^{\mathrm{ref}}_i,y;\ell).
\end{equation}
The preference loss is
\begin{equation}
  \mathcal{L}_{\mathrm{DPO}}
  = -\log \sigma\left(
    \beta\left[
    (s_\theta(y^+) - s_\theta(y^-))
    - (s_{\mathrm{ref}}(y^+) - s_{\mathrm{ref}}(y^-))
    \right]\right),
\end{equation}
and the reported preference stage uses
\begin{equation}
  \mathcal{L}_{\mathrm{pref}}
  = \mathcal{L}_{\mathrm{DPO}}
  + \lambda_{\mathrm{sft}}\min_{i \le M}d_\delta(\hat{y}_i,y^+;\ell).
\end{equation}
This final stage sharpens MixiMotion without asking the preference loss to
learn the full motion prior from scratch.

\subsection{Training Algorithm}

\noindent\textbf{Algorithm 1: Progressive distillation for MixiMotion.}
\begin{enumerate}
\item Build a curated prompt bank from text-to-motion descriptions.
\item For each prompt $p$, run the frozen HY-Motion-1.0-Lite teacher with $K$
deterministic seeds and cache $\mathcal{Y}_T(p)$, text features, and motion
latents.
\item Normalize cached and SSAE prompt text with the same deterministic rule,
remove every exact match from the teacher cache, and verify that the remaining
exact normalized-text overlap is zero.
\item Train MixiMotion with $\mathcal{L}_{\mathrm{base}}$ using $M$ randomly
sampled noise candidates per prompt and bidirectional latent/body-space
matching.
\item Continue with gait/body matching, then oversample locomotion prompts by
$6\times$ for the hard-locomotion subphase.
\item Continue with the semantic-balanced curriculum, oversampling locomotion
and stationary/localized prompts by $4\times$ and applying group-specific gait
or root-stability terms.
\item Construct preference pairs $(p,y^+,y^-)$ from the cached teacher
candidates using the group-aware score $q(y,p)$.
\item Freeze a reference copy of the pre-preference student, then optimize
$\mathcal{L}_{\mathrm{pref}}$ with the same noise tokens for the student and
reference candidates.
\item Return the exponential-moving-average student as MixiMotion.
\end{enumerate}

\begin{table}
\caption{Selected implementation details for the reported MixiMotion training
run.}
\label{tab:training-details}
\centering
\begin{tabular}{ll}
\hline
Item & Value \\
\hline
Teacher for cached samples & HY-Motion-1.0-Lite \\
Teacher-bank rows after filtering & 48,065 \\
Exact normalized-text SSAE overlap & 0 \\
Teacher candidates per prompt $K$ & 4 \\
Student candidates, base distillation $M$ & 8 \\
Student candidates, preference stage $M$ & 4 \\
Motion length & 120 frames \\
Optimizer & AdamW \\
Base stage & 18 epochs at $8.0 \times 10^{-5}$ \\
Gait/body stage & 10 epochs at $3.0 \times 10^{-5}$ \\
Hard-locomotion stage & 6 epochs at $2.0 \times 10^{-5}$ \\
Semantic-balance stage & 4 epochs at $1.0 \times 10^{-5}$ \\
Preference stage & 3 epochs at $2.0 \times 10^{-6}$ \\
Batch size & 4 \\
DPO inverse temperature $\beta$ & 8.0 \\
SFT weight $\lambda_{\mathrm{sft}}$ & 0.35 \\
Smooth-L1 transition $\delta$ & 0.35 \\
Network evaluations per candidate & 1 \\
EMA decay & 0.999 \\
Base-stage split & 70\% / 15\% / 15\% (train/val/test) \\
Preference-stage split & 90\% / 10\% (train/val) \\
Precision & bfloat16 \\
\hline
\end{tabular}
\end{table}

Figure~\ref{fig:curriculum} illustrates the curriculum as four consecutive
training stages rather than as implementation-version checkpoints.

\begin{figure}[t]
\centering
\includegraphics[width=\textwidth]{figures/miximotion_curriculum.png}
\caption{Progressive curriculum used to train MixiMotion. Cached teacher
supervision is reused across all stages; solid arrows denote EMA checkpoint
handoff. The sequence does not imply monotonic improvement.}
\label{fig:curriculum}
\end{figure}

\section{Experimental Setup}

\subsection{Training Data}

The reported model uses a curated prompt bank collected from
HY-Motion-1.0 examples, ViMoGen motion annotations~\cite{vimogen2026}, and manually filtered
text-to-motion descriptions. For each prompt, HY-Motion-1.0-Lite generates four
teacher candidates at a fixed length of 120 frames. Before the reported
retraining run, we normalize text by stripping outer whitespace, converting to
lowercase, and collapsing repeated whitespace. We remove all 1,935 cached rows
whose normalized text matches an SSAE prompt; these rows cover 1,883 unique
SSAE prompts. The resulting teacher bank contains 48,065 samples, and a
post-filter audit finds zero exact normalized-text overlap with all 1,969 SSAE
prompts. All curriculum stages and the preference bank are then rebuilt from
this filtered cache. This audit excludes exact text matches but does not rule
out semantic paraphrases or near-duplicate descriptions.

\subsection{Evaluation Protocol}

We evaluate on SSAE~\cite{hymotion2025}, which decomposes each text prompt into
fine-grained yes/no questions about visible motion semantics. We render one
video per SSAE index and per generator, producing 1,969 videos for each system
and 7,876 in total. A VLM answers all questions for each video, and the score
is the ratio of yes answers. Our native-video protocol passes all 120 frames of
each rendered sequence to Qwen/Qwen3-VL-30B-A3B-Instruct-FP8~\cite{qwen3vl2025};
no temporal frame subsampling is used.

The primary comparison uses two motion generators:
\begin{itemize}
\item \textbf{MixiMotion}: our current one-step HY-Motion-1.0 model retrained
from the exact-text-filtered teacher bank.
\item \textbf{HY-Motion-1.0 full}: the public multi-step HY-Motion-1.0 model
rendered with the same SSAE prompts and 120-frame length. Each benchmark run
uses the deterministic seed recorded in its manifest.
\end{itemize}
For external context, we also run SSAE on two public text-to-motion baselines
with available checkpoints: \textbf{MotionHiFlow}~\cite{motionhiflow2026} and
\textbf{MotionLCM}~\cite{motionlcm2024}. All systems use the same prompt-index
mapping, 120-frame output length, SSAE questions, and native-video judge
settings. MixiMotion and HY-Motion-1.0 full share the same 720$\times$720
skeleton renderer, camera, and 30 FPS encoding (4 s). The public baselines use
their native project renderers at 400$\times$400 and 300$\times$300,
respectively, and are encoded at 20 FPS (6 s). We do not temporally rescale
these motions because doing so changes their native execution speed;
consequently, public-baseline scores are contextual rather than strictly
renderer- and duration-controlled comparisons.

For VLM judging, we use one all-frame native-video protocol:
\begin{itemize}
\item \textbf{Qwen/Qwen3-VL-30B-A3B-Instruct-FP8}, used for native-video
evaluation with all 120 frames, identical SSAE questions, temperature zero,
a 256-token response limit, and a strict yes/no JSON schema for every motion
generator.
\end{itemize}
Strict structured decoding succeeded directly for 7,873 of the 7,876 videos.
For three videos, the backend entered a whitespace loop after opening the JSON
object. We reissued each question independently to the same model with the
same 120-frame video, temperature, and yes/no instruction, changing only the
output-format constraint. All four systems consequently contain 1,969
successful videos and the same 6,418 answered questions.

\subsection{Human Evaluation Protocol}

To complement the VLM-based assessment, we conduct a blinded human evaluation
with 25 participants familiar with computer graphics or motion synthesis. The
study covers 300 SSAE prompts per system. For each prompt, model outputs are
shown in randomized pairs with model identities and left--right ordering
hidden. Participants assign an independent five-point Likert score to each
video for: (1) semantic alignment with the input prompt, (2) motion
naturalness, and (3) overall motion quality, including foot--ground contact,
stability, and visible artifacts. The reported values aggregate ratings over
the 300 prompts for each system. We report the mean and standard error, and use
Krippendorff's alpha to summarize inter-rater agreement.

\section{Results}

\subsection{Curriculum Ablation}

We use a cumulative ablation from the preceding development run to examine the
effect of each curriculum stage. All variants use the same architecture,
teacher cache, rendering pipeline, and native-video SSAE judge; only the
indicated curriculum component is added. In addition to SSAE, we report
foot-skating and contact metrics for every checkpoint. These exploratory
ablation checkpoints predate exact-text filtering and are not mixed with the
clean-retrained checkpoint in the main comparison.

\begin{table}
\caption{Exploratory cumulative ablation from the preceding development run.
Higher SSAE and contact scores are better; lower foot skating is better.
Physical metrics are dataset means over the 1,969 development motions. Best
values are in bold.}
\label{tab:curriculum-ablation}
\centering
\small
\setlength{\tabcolsep}{6pt}
\renewcommand{\arraystretch}{1.2}
\begin{tabular*}{\textwidth}{@{\extracolsep{\fill}}lccc@{}}
\hline
Variant & \shortstack{Native-video\\SSAE} & Foot skating & Contact score \\
\hline
Base distillation & 0.8822 & \textbf{0.0401} & \textbf{0.7901} \\
$+$ Motion-quality & 0.8802 & 0.0495 & 0.7621 \\
$+$ Semantic balance & \textbf{0.8824} & 0.0589 & 0.7376 \\
$+$ Preference/DPO (development) & 0.8813 & 0.0558 & 0.7513 \\
\hline
\end{tabular*}
\end{table}

Base distillation and semantic balancing obtain nearly identical SSAE scores.
The motion-quality and final development preference checkpoints are lower than the base by
0.0020 and 0.0009, respectively. The aggregate physical metrics are also
non-monotonic: later stages have more foot skating and less contact than the
base, while preference optimization partly reverses the changes introduced by
semantic balancing. Thus, this ablation does not support a claim that every
curriculum stage improves these aggregate metrics. For the development DPO
checkpoint, the foot-skating ratio is $0.0558\pm0.1433$ and the contact ratio
is $0.7513\pm0.3347$ across 1,969 motions. The same renderer-side implementation
and thresholds are used for every checkpoint.

\subsection{Model Size and Inference Cost}

Table~\ref{tab:cost-results} compares model size, solver cost, and measured
inference characteristics. Parameter counts are computed from saved generation
checkpoints and exclude text encoders, VLM judges, and video renderers. The two
HY-Motion-family rows count the motion transformer. MotionHiFlow combines its
flow generator and VAE, whereas MotionLCM combines its denoiser and VAE; these
external parameter scopes are therefore contextual rather than
architecture-identical. Every system is profiled on the same RTX 5090 with
batch size 1, bfloat16 precision, and 120 output frames. We perform 20
unreported warm-up iterations followed by 100 timed iterations and synchronize
CUDA before and after every measurement. Generation-only latency covers the
native motion generator/solver. Model-native end-to-end latency additionally
includes text encoding and each project's motion decoder, but excludes video
rendering and VLM judging. HY-Motion-1.0 full and MixiMotion terminate at the
same decoded HY-Motion representation. MotionHiFlow terminates at decoded
motion features, while MotionLCM additionally converts features to joints and
copies them to CPU. Thus, common hardware and timing controls are shared, but
external end-to-end endpoints are not identical. We report mean, standard
deviation, and 95th percentile. Peak VRAM is measured after resetting the CUDA
peak-memory counter immediately before each timed scope.

\begin{table}
\caption{Controlled model-size and inference-cost profiling on 120-frame
samples. Latency and memory use one RTX 5090/BF16 harness at batch size 1.
External model-native endpoints differ as described in the text. Lower is
better for every numeric column; best values are in bold.}
\label{tab:cost-results}
\centering
\small
\setlength{\tabcolsep}{6pt}
\renewcommand{\arraystretch}{1.18}
\textit{(a) Model complexity and memory}\\[2pt]
\begin{tabular*}{\textwidth}{@{\extracolsep{\fill}}lccc@{}}
\hline
Model &
\shortstack{Generator\\params (M)} &
\shortstack{Sampling\\steps/NFE} &
\shortstack{Peak VRAM\\(GiB)} \\
\hline
HY-Motion-1.0 full & 1042.9 & 50 & 23.23 \\
MixiMotion (ours) & 460.0 & \textbf{1} & 20.73 \\
MotionHiFlow & 33.1 & 12 & \textbf{0.43} \\
MotionLCM & \textbf{26.5} & \textbf{1} & 1.45 \\
\hline
\end{tabular*}

\vspace{6pt}
\textit{(b) Latency}\\[2pt]
\setlength{\tabcolsep}{4pt}
\begin{tabular*}{\textwidth}{@{\extracolsep{\fill}}lcccc@{}}
\hline
Model &
\shortstack{Generation\\mean $\pm$ std (ms)} &
\shortstack{Generation\\P95 (ms)} &
\shortstack{End-to-end\\mean $\pm$ std (ms)} &
\shortstack{End-to-end\\P95 (ms)} \\
\hline
HY-Motion-1.0 full & $829.58\pm3.19$ & 834.68 & $968.25\pm3.69$ & 974.94 \\
MixiMotion (ours) & $9.30\pm0.39$ & 9.94 & $143.55\pm1.29$ & 145.67 \\
MotionHiFlow & $111.19\pm15.98$ & 139.02 & $119.68\pm15.89$ & 148.68 \\
MotionLCM & $\mathbf{3.04\pm0.24}$ & \textbf{3.35} & $\mathbf{13.94\pm0.47}$ & \textbf{14.85} \\
\hline
\end{tabular*}
\end{table}

MixiMotion uses 44.1\% of the HY-Motion-1.0 full motion-generator parameters
and one network evaluation rather than 50 solver evaluations. Under the same
HY-Motion-family scope, it is $89.2\times$ faster for generation alone and
$6.75\times$ faster end to end. The smaller end-to-end gain reflects shared
text-encoding and motion-decoding costs. In the model-native profiles,
MixiMotion's one-step generator is $11.9\times$ faster than MotionHiFlow's
12-step solver, and MotionLCM records the lowest latency and parameter count.
Because their decoded endpoints differ, these cross-family latency values
provide deployment context rather than a strict pipeline-level speedup claim.
MotionLCM's lower SSAE score in Table~\ref{tab:native-video-ssae} further shows
that computational cost and semantic alignment should be considered together.

\subsection{Native-Video SSAE Evaluation}

We evaluate all four generators with
Qwen/Qwen3-VL-30B-A3B-Instruct-FP8 through its native video interface. Each
request contains all 120 frames of one motion video and the same SSAE yes/no
questions used for every system. The model identifier, frame count,
temperature, response schema, and parsing procedure are held fixed. This
protocol replaces the earlier frame-sampled Qwen experiment.

\begin{table}
\caption{SSAE scores under the all-frame native-video
Qwen/Qwen3-VL-30B-A3B-Instruct-FP8 judge. MixiMotion is retrained after
removing exact normalized-text SSAE matches from its teacher bank. Every system
has 1,969 successful videos and 6,418 answered questions. Higher is better;
the best score in each row is in bold.}
\label{tab:native-video-ssae}
\centering
\small
\setlength{\tabcolsep}{4pt}
\renewcommand{\arraystretch}{1.15}
\begin{tabular*}{\textwidth}{@{\extracolsep{\fill}}lcccc@{}}
\hline
Category &
\shortstack{MixiMotion\\(ours)} &
\shortstack{HY-Motion-1.0\\full} &
\shortstack{Motion\\HiFlow} &
\shortstack{Motion\\LCM} \\
\hline
Daily      & 0.9009 & \textbf{0.9075} & 0.8617 & 0.8065 \\
Fitness    & 0.7487 & 0.7613 & \textbf{0.7717} & 0.7213 \\
Game       & 0.9205 & \textbf{0.9215} & 0.9187 & 0.8977 \\
Locomotion & \textbf{0.9757} & \textbf{0.9757} & 0.9689 & 0.9621 \\
Social     & 0.9332 & 0.9419 & \textbf{0.9452} & 0.9211 \\
Sports     & 0.8430 & \textbf{0.8451} & 0.8274 & 0.7807 \\
\hline
\textbf{Overall} & 0.8802 & \textbf{0.8856} &
0.8764 & 0.8415 \\
\hline
\end{tabular*}
\end{table}

MixiMotion retains 99.4\% of the HY-Motion-1.0 full score, with an absolute
gap of 0.0055. It matches the full model on locomotion at the reported
precision and is within 0.0021 on game and sports prompts. MotionHiFlow is
strongest on fitness and social among the compact models, while MixiMotion has
the numerically highest compact-model overall score. These small differences should be
interpreted cautiously because the exact-text audit does not exclude semantic
near-duplicates and the experiment uses a single deterministic VLM judge. A
paired video-level bootstrap with 10,000 resamples gives a 95\% confidence
interval of $[0.8687, 0.8915]$ for MixiMotion and $[0.8746, 0.8964]$ for
HY-Motion-1.0 full. The paired difference (MixiMotion minus full) is $-0.0055$
with a 95\% interval of $[-0.0121, 0.0011]$, which includes zero.
MotionHiFlow's paired difference is $-0.0092$ with interval
$[-0.0186, 0.0002]$, whereas MotionLCM's difference is $-0.0441$ with interval
$[-0.0553, -0.0332]$. Thus, under this protocol, the VLM scores do not resolve
a statistically clear difference between MixiMotion and the full baseline.

\subsection{Human Evaluation}

Table~\ref{tab:human-evaluation} reports the human ratings. HY-Motion-1.0 full
has the highest mean score on all three criteria. MixiMotion is the strongest
compact system, exceeding MotionHiFlow and MotionLCM on semantic alignment,
naturalness, and overall motion quality. Its descriptive gap from the full
baseline is 0.16 for semantic alignment and 0.17 for both naturalness and
overall quality. The aggregate inter-rater agreement is
Krippendorff's $\alpha=0.68$.

\begin{table}[H]
\caption{Blinded human evaluation on a five-point Likert scale. Values are
means over 300 prompts $\pm$ standard error. Higher is better. Bold denotes
the best overall result and underlining denotes the best compact model.}
\label{tab:human-evaluation}
\centering
\small
\setlength{\tabcolsep}{4pt}
\renewcommand{\arraystretch}{1.12}
\begin{tabular*}{\textwidth}{@{\extracolsep{\fill}}lccc@{}}
\hline
Model &
\shortstack{Semantic\\alignment} &
Naturalness &
\shortstack{Overall motion\\quality} \\
\hline
HY-Motion-1.0 full & $\mathbf{4.48\pm0.04}$ & $\mathbf{4.52\pm0.03}$ & $\mathbf{4.50\pm0.04}$ \\
MixiMotion (ours) & $\underline{4.32\pm0.05}$ & $\underline{4.35\pm0.04}$ & $\underline{4.33\pm0.05}$ \\
MotionHiFlow & $4.18\pm0.06$ & $4.22\pm0.05$ & $4.20\pm0.06$ \\
MotionLCM & $3.95\pm0.07$ & $4.01\pm0.06$ & $3.98\pm0.07$ \\
\hline
\end{tabular*}
\end{table}

These descriptive results indicate that MixiMotion preserves much of the
perceptual quality of the full HY-Motion-1.0 baseline while improving over the
evaluated efficient baselines. They complement rather than replace the
automated semantic and physical measurements.

Figure~\ref{fig:qualitative} shows representative deterministic outputs from
the two systems in the HY-Motion family. We render saved SMPL parameters directly
with the same Three.js wooden-mesh visualizer used by the Streamlit demo. This
comparison is intended to expose temporal pose differences rather than replace
quantitative or human evaluation.

\begin{figure}[t]
\centering
\includegraphics[width=\textwidth]{figures/qualitative_comparison.png}
\caption{Qualitative comparison on four SSAE prompts using the Streamlit
Three.js wooden-mesh renderer. Each row shows three temporally ordered
keyframes from the corresponding 120-frame output. For locomotion, early
frames are selected because the demo camera remains fixed while the body
translates. Both columns use identical rendering code and camera
initialization.}
\label{fig:qualitative}
\end{figure}

\section{Discussion}

\subsection{Why a Curriculum Instead of a Single Stage?}

The preference loss does not by itself teach the full motion prior, so the base
distillation remains necessary to initialize the one-step mapping. The later
stages target gait/contact behavior, semantic balance, and candidate
preference. However, Table~\ref{tab:curriculum-ablation} shows that their
effects are not captured as monotonic gains in aggregate native-video SSAE,
foot skating, or contact ratio. More targeted per-group physical metrics and
an independently sourced prompt benchmark are needed to establish which later
stages are beneficial.

\subsection{Scaling the Development Data}

Daily and social prompts contain many fine-grained hand, torso, and interaction
cues that are difficult for compact one-step generators. We therefore expect
larger, curated prompt banks to be particularly useful for these behaviors.
Future scaling should use held-out validation prompts and human inspection
rather than selecting examples from benchmark failures.

\subsection{Limitations}

First, MixiMotion is distilled from HY-Motion-1.0-Lite rather than
from the full HY-Motion-1.0 teacher. This may impose a ceiling on semantic
alignment. Second, although exact normalized-text SSAE matches are removed
before retraining, we have not performed semantic near-duplicate detection
across the remaining training descriptions; the result therefore establishes
exact-text decontamination rather than complete semantic independence. Third,
SSAE uses VLM judgments over rendered skeleton videos, so the resulting scores
remain dependent on the selected judge. Although the human study provides an
independent perceptual assessment, it is limited to 25 participants and 300
prompts per system. SSAE also measures semantic question answering rather than
distributional fidelity or motion diversity. Fourth, MixiMotion is evaluated
at a fixed 120-frame length. The HY-Motion-family
videos are encoded at 30 FPS, whereas the public baselines are encoded at
20 FPS, so the baseline comparison is not perfectly duration controlled.
Variable-length and prompt-engineered generation remain future work.

\section{Conclusion}

We presented MixiMotion, our one-step HY-Motion-1.0 model trained with a
progressive curriculum distillation recipe for text-to-motion generation.
MixiMotion uses one-step inference and 460.0M rather than 1.04B
motion-generator parameters. After retraining on a teacher bank with zero
exact normalized-text SSAE overlap, MixiMotion obtains 0.8802 SSAE versus
0.8856 for HY-Motion-1.0 full in the all-frame native-video evaluation,
retaining 99.4\% of its score and obtaining the highest numerical score among
the evaluated compact models. Public-baseline comparisons remain contextual
because their native renderers and frame rates differ. The method provides a
practical route toward interactive text-to-motion systems. The blinded human
study likewise places MixiMotion above the evaluated compact baselines and
slightly below HY-Motion-1.0 full, while evaluation on an independently
sourced benchmark remains important for establishing broader generalization.

\begin{credits}
\subsubsection{\ackname} The authors thank the maintainers of HY-Motion-1.0,
SSAE, ViMoGen, and the open-source text-to-motion ecosystem for releasing
models, prompts, and evaluation tools.

\subsubsection{\discintname}
The authors have no competing interests to declare that are relevant to the
content of this article.
\end{credits}

\begin{thebibliography}{99}

\bibitem{hymotion2025}
Tencent Hunyuan 3D Digital Human Team: HY-Motion-1.0: Scaling Flow Matching
Models for Text-to-Motion Generation. arXiv preprint arXiv:2512.23464 (2025)

\bibitem{moflow2024}
Wang, Y., et al.: MoFlow: One-Step Flow Matching for Human Motion Generation.
arXiv preprint arXiv:2503.09950 (2025)

\bibitem{motionlcm2024}
Chen, L., et al.: MotionLCM: Real-time Controllable Motion Generation via
Latent Consistency Model. arXiv preprint arXiv:2404.19759 (2024)

\bibitem{motionhiflow2026}
Li, H., Lin, X., Zeng, L.-A., Kang, Y., Li, S., Hu, J.-F.:
MotionHiFlow: Text-to-Motion via Hierarchical Flow Matching. In:
Proceedings of the IEEE/CVF Conference on Computer Vision and Pattern
Recognition (CVPR) (2026)

\bibitem{mdm2022}
Tevet, G., Raab, S., Gordon, B., Shafir, Y., Cohen-Or, D., Bermano, A.H.:
Human Motion Diffusion Model. In: ICLR (2023)

\bibitem{humanml3d2022}
Guo, C., Zou, S., Zuo, X., Wang, S., Ji, W., Li, X., Cheng, L.:
Generating Diverse and Natural 3D Human Motions from Text. In: CVPR (2022)

\bibitem{t2mgpt2023}
Zhang, J., Zhang, Y., Cun, X., Huang, S., Zhang, Y., Zhao, H., Lu, H.,
Shan, Y.: Generating Human Motion from Textual Descriptions with Discrete
Representations. In: CVPR (2023)

\bibitem{mld2023}
Chen, X., Jiang, B., Liu, W., Huang, Z., Fu, B., Chen, T., Yu, G.:
Executing Your Commands via Motion Diffusion in Latent Space. In: CVPR (2023)

\bibitem{vimogen2026}
ViMoGen: Large-scale video-to-motion prompt and motion dataset. Hugging Face
dataset repository, \url{https://huggingface.co/datasets/wruisi/ViMoGen-228K}

\bibitem{progressivedistill2022}
Salimans, T., Ho, J.: Progressive Distillation for Fast Sampling of Diffusion
Models. In: International Conference on Learning Representations (ICLR) (2022)

\bibitem{consistencymodels2023}
Song, Y., Dhariwal, P., Chen, M., Sutskever, I.: Consistency Models. In:
Proceedings of the 40th International Conference on Machine Learning (ICML),
pp. 32211--32252 (2023)

\bibitem{imle2018}
Li, K., Malik, J.: Implicit Maximum Likelihood Estimation. arXiv preprint
arXiv:1809.09087 (2018)

\bibitem{dpo2023}
Rafailov, R., Sharma, A., Mitchell, E., Manning, C.D., Ermon, S., Finn, C.:
Direct Preference Optimization: Your Language Model Is Secretly a Reward
Model. In: Advances in Neural Information Processing Systems 36 (2023)

\bibitem{qwen3vl2025}
Bai, S., et al.: Qwen3-VL Technical Report. arXiv preprint arXiv:2511.21631
(2025)

\end{thebibliography}

\end{document}
